\documentclass[11pt]{article}
\usepackage[a4paper,margin=28mm]{geometry}
\usepackage[T1]{fontenc}
\usepackage{lmodern}
\usepackage{amsmath,amssymb,amsthm,booktabs,array}
\usepackage[hidelinks]{hyperref}
\usepackage{microtype}
\newtheorem{theorem}{Theorem}
\newtheorem{lemma}[theorem]{Lemma}
\newtheorem{proposition}[theorem]{Proposition}
\newtheorem{remark}[theorem]{Remark}
\newcommand{\R}{\mathbb R}
\newcommand{\N}{\mathbb N}
\title{Sharp root log-concavity of alternating Motzkin sums}
\author{Dev D. Goyal}
\date{26 October 2026}
\usepackage[numbers]{natbib}
\usepackage{xurl}
\begin{document}
\maketitle

\begin{abstract}
Let $M_n$ be the Motzkin numbers and let
$T_n=\sum_{j=0}^n(-1)^{n-j}M_j$.
The sequence $T_n^{1/n}$ is strictly log-concave on its tail
starting at $n=5$. The inequality fails at center $n=5$, so no earlier tail
is log-concave. This proves Conjecture~1 of Zhao, which asks for
log-concavity from index nine. The argument combines established ratio
bounds with a size estimate. These estimates prove the infinite range;
exact integer comparisons cover the remaining 25 centers. A Lean
development proves the real-root inequality and its sharpness. An added
source-definition module identifies the recursive implementation with
the finite sum used in the conjecture.
\end{abstract}

\noindent\textbf{Verification.}
The source theorem uses only \texttt{propext}, \texttt{Classical.choice},
and \texttt{Quot.sound}. The repository's
\href{https://github.com/dvgyl/math-and-theory/blob/main/verification/summary.md}{verification record}
gives the exact builds, statement comparisons, independent kernel checks,
and AI review status.

\section{The problem and the result}

A positive sequence $(a_n)$ is log-concave if
$a_n^2\ge a_{n-1}a_{n+1}$ at each interior index.
Strict log-concavity means that every such inequality is strict.
For a tail starting at $s$, the interior indices are $n\ge s+1$.
This distinction between a tail index and a center index is needed below.

The Motzkin numbers satisfy the source recurrence and initial values:
\begin{equation}\label{eq:motzkin}
M_0=M_1=1,\qquad
(n+3)M_{n+1}=(2n+3)M_n+3nM_{n-1}\quad(n\ge1).
\end{equation}
These conditions determine the sequence uniquely over $\mathbb Q$ or $\mathbb R$.
Define
\begin{equation}\label{eq:alternating}
T_n=\sum_{j=0}^{n}(-1)^{n-j}M_j.
\end{equation}
The first values are
\[
T_0,T_1,\ldots,T_9=1,0,2,2,7,14,37,90,233,602.
\]
Cancellation of consecutive sums gives
\begin{equation}\label{eq:Trec}
T_{n+1}=M_{n+1}-T_n.
\end{equation}

Zhao studies this sequence in \cite{Zhao}. His Conjecture~1 asks whether
$(T_n^{1/n})_{n\ge9}$ is log-concave. The following theorem gives a sharper statement.

\begin{theorem}\label{thm:main}
The values $T_n$ are positive for $n\ge2$. For every integer $n\ge6$,
\begin{equation}\label{eq:main}
T_n^{\,2(n^2-1)}>
T_{n-1}^{\,n(n+1)}T_{n+1}^{\,n(n-1)}.
\end{equation}
Thus $(T_n^{1/n})_{n\ge5}$ is strictly log-concave.
At center five the inequality is strictly reversed:
\begin{equation}\label{eq:failure}
(T_5^{1/5})^2<T_4^{1/4}T_6^{1/6}.
\end{equation}
Consequently, five is the smallest starting index of a log-concave tail.
\end{theorem}

The ratio bounds used in the proof are established ingredients.
The Motzkin bounds come from the interlacing method of Do\v{s}li\'c and
Veljan \cite{DV}. Zhao proves an interlacing estimate for the alternating
sums \cite[eqs.~(7)--(8)]{Zhao}. Proofs of the needed estimates are included.
General methods for root log-concavity also appear in
Chen, Guo, and Wang \cite{CGW} and in Hou and Li \cite{HL}.
The contribution here is the stated alternating-sum result and its sharp
threshold. No claim is made that the general method is new.

The proof has two parts. First, the proof bounds the logarithm of the ratio
between the two sides of \eqref{eq:main}. The bound is positive for
$n\ge31$. Second, the proof compares exact integers for $6\le n\le30$.
This covers every required center.

\section{Ratio and size estimates}

Put
\[
x_n=\frac{M_{n+1}}{M_n},\qquad
b_n=\frac{6(n+1)}{2n+5}.
\]
The recurrence shows by induction that $M_n>0$ for every $n$.

\begin{lemma}\label{lem:ratios}
For every integer $n\ge1$,
\begin{equation}\label{eq:ratios}
b_n\le x_n\le b_{n+1}.
\end{equation}
Moreover, $0<T_n\le M_n$ for $n\ge2$.
\end{lemma}

\begin{proof}
At $n=1$, the values satisfy $x_1=2=b_2$ and $b_1\le2$.
For $n\ge2$, define
\[
f_n(t)=\frac{2n+3}{n+3}+\frac{3n}{(n+3)t}\quad(t>0).
\]
The recurrence gives $x_n=f_n(x_{n-1})$.
The function $f_n$ is decreasing. Thus the induction hypothesis gives
\[
f_n(b_n)\le x_n\le f_n(b_{n-1}).
\]
Direct calculation yields
\[
f_n(b_n)-b_n
=\frac{3(n-2)}{2(n+1)(n+3)(2n+5)}\ge0,
\]
and
\[
b_{n+1}-f_n(b_{n-1})
=\frac{9}{2(n+3)(2n+7)}>0.
\]
This proves \eqref{eq:ratios}.
Since $b_n>1$ for $n\ge1$, the Motzkin numbers increase strictly from
index one. Now $T_1=0$ and $T_2=2=M_2$.
If $0\le T_{n-1}\le M_{n-1}$ and $n\ge2$, then
\[
M_n-M_{n-1}\le T_n=M_n-T_{n-1}\le M_n.
\]
The lower bound is positive. This proves the remaining claim by induction.
\end{proof}

\begin{lemma}\label{lem:interlace}
Let $y_n=T_{n+1}/T_n$ for $n\ge2$. Then
\begin{equation}\label{eq:interlace}
x_n\le y_n\le x_{n+1}\qquad(n\ge9).
\end{equation}
\end{lemma}

\begin{proof}
Equation~\eqref{eq:Trec} gives
\begin{equation}\label{eq:yrec}
y_{n+1}=x_{n+1}-1+\frac{x_{n+1}}{y_n}.
\end{equation}
At $n=9$, the values are
\[
x_9=\frac{2188}{835},\quad y_9=\frac{793}{301},\quad
x_{10}=\frac{2899}{1094}.
\]
The two differences are $3567/251335$ and $5057/329294$, respectively.
Both are positive.

Suppose \eqref{eq:interlace} holds at $n\ge9$.
Its upper bound and \eqref{eq:yrec} imply $y_{n+1}\ge x_{n+1}$.
For the other bound, put
\[
f(t)=\frac{2n+5}{n+4}+\frac{3(n+1)}{(n+4)t},\quad
c=\frac{2n+7}{n+5},\quad d=\frac{3(n+2)}{n+5},
\]
and define
\[
H(t)=f(t)-1+\frac{f(t)}t-c-\frac d{f(t)}\quad(t>0).
\]
The functions $f(t)$, $f(t)/t$, and $-d/f(t)$ are decreasing.
For the middle function, write it as $A/t+B/t^2$ with $A,B>0$.
Thus $H$ is decreasing. The recurrence and the induction hypothesis give
\begin{align*}
y_{n+1}-x_{n+2}
&\le x_{n+1}-1+\frac{x_{n+1}}{x_n}-x_{n+2}\\
&=H(x_n)\le H(b_n)\\
&=-\frac{3(4n^2-13n-53)}{4(n+1)(n+4)(n+5)(2n+5)}<0.
\end{align*}
The numerator is positive for $n\ge9$.
This proves the induction step.
\end{proof}

\begin{lemma}\label{lem:size}
For every integer $n\ge0$,
\begin{equation}\label{eq:size}
M_n\le\frac{3^n}{n+1}.
\end{equation}
\end{lemma}

\begin{proof}
The cases $n=0,1$ hold directly. For $n\ge1$, apply the recurrence and
the bounds at $n$ and $n-1$ to obtain
\[
M_{n+1}\le\frac{3^n(3n+4)}{(n+1)(n+3)}
<\frac{3^{n+1}}{n+2}.
\]
The strict comparison follows from
$3(n+1)(n+3)-(3n+4)(n+2)=2n+1>0$.
\end{proof}

\section{The infinite range}

For $n\ge3$, put $q_n=T_n/T_{n-1}$ and define
\[
D_n=n(n^2-1)
\left(\frac{2\log T_n}{n}
-\frac{\log T_{n-1}}{n-1}
-\frac{\log T_{n+1}}{n+1}\right).
\]
All values in the logarithms are positive by Lemma~\ref{lem:ratios}.
The factor $n(n^2-1)$ is positive. Hence $D_n>0$ is equivalent both to
strict root log-concavity at center $n$ and to \eqref{eq:main}.
Substitution gives the useful identity
\begin{equation}\label{eq:gap}
D_n=2n\log q_n
-n(n-1)\log\frac{q_{n+1}}{q_n}-2\log T_n.
\end{equation}

For $n\ge10$, Lemmas~\ref{lem:ratios} and \ref{lem:interlace} give
\[
q_n\ge b_{n-1}=\frac{3n}{n+3/2},\qquad
\frac{q_{n+1}}{q_n}
\le\frac{b_{n+2}}{b_{n-1}}
=1+\frac9{n(2n+9)}.
\]
Also, Lemmas~\ref{lem:ratios} and \ref{lem:size} imply
\[
\log T_n\le n\log3-\log(n+1).
\]
Use these estimates in \eqref{eq:gap}. Since $\log(1+t)<t$ for $t>0$,
\begin{align}
D_n
&\ge 2\log(n+1)
-2n\log\left(1+\frac3{2n}\right)
-n(n-1)\log\left(1+\frac9{n(2n+9)}\right)\notag\\
&>2\log(n+1)-3-\frac{9(n-1)}{2n+9}
=:F(n).\label{eq:F}
\end{align}

For real $x\ge10$,
\[
F'(x)=\frac{8x^2-27x+63}{(x+1)(2x+9)^2}>0.
\]
To give an exact positive value at the threshold, observe that
\[
\log2=2\int_0^{1/3}\frac{dt}{1-t^2}
>2\int_0^{1/3}(1+t^2)\,dt
=\frac{56}{81}.
\]
The strict inequality follows because
$1/(1-t^2)>1+t^2$ for $0<t<1/3$.
Consequently,
\[
F(31)=10\log2-3-\frac{270}{71}
>\frac{560}{81}-3-\frac{270}{71}
=\frac{637}{5751}>0.
\]
Monotonicity of $F$ and \eqref{eq:F} prove \eqref{eq:main} for every
integer $n\ge31$.

\section{The finite range and the sharp threshold}

It remains to check the 25 centers $6\le n\le30$.
Table~\ref{tab:values} gives the values required for those comparisons,
together with the value $T_4$ needed for sharpness.
Every value follows from \eqref{eq:motzkin} and \eqref{eq:Trec}.
Substitution in \eqref{eq:main} gives a strict inequality at each of the
25 centers. These comparisons use integers only.

\begin{table}[ht]
\centering\small
\begin{tabular}{rr@{\qquad}rr}
\toprule
$n$&$T_n$&$n$&$T_n$\\\midrule
4&7&18&4805596\\
5&14&19&13393688\\
6&37&20&37458331\\
7&90&21&105089228\\
8&233&22&295673995\\
9&602&23&834086420\\
10&1586&24&2358641377\\
11&4212&25&6684761124\\
12&11299&26&18985057352\\
13&30536&27&54022715450\\
14&83098&28&154000562759\\
15&227474&29&439742222070\\
16&625993&30&1257643249141\\
17&1730786&31&3602118427250\\\bottomrule
\end{tabular}
\caption{Exact alternating Motzkin sums.}\label{tab:values}
\end{table}

The following complete arithmetic procedure reproduces the finite
comparisons. The division in the recurrence is checked to be exact.
\begin{verbatim}
M = [1, 1]
for n in range(1, 31):
    numerator = (2*n + 3)*M[n] + 3*n*M[n-1]
    assert numerator % (n + 3) == 0
    M.append(numerator // (n + 3))
T = [sum((-1)**(n-j)*M[j] for j in range(n+1))
     for n in range(32)]
for n in range(6, 31):
    assert T[n]**(2*(n*n-1)) > (
        T[n-1]**(n*(n+1))*T[n+1]**(n*(n-1)))
assert T[5]**48 < T[4]**30*T[6]**20
\end{verbatim}

At center five, the reverse comparison follows from
\[
37^{10}-2^{24}7^9=4131563191399737>0.
\]
Indeed, $T_4=7$, $T_5=14$, and $T_6=37$.
Squaring the inequality and multiplying by $7^{30}$ gives
$14^{48}<7^{30}37^{20}$.
It is equivalent to \eqref{eq:failure}.
Every tail starting at an index at most four contains center five.
The infinite argument and the finite comparisons prove
Theorem~\ref{thm:main}.

\section{Formal statements and reproducibility}

The analytic Lean development defines rational Motzkin numbers through
adjacent pairs. It defines the alternating sequence recursively through
\eqref{eq:Trec}. A separate source-definition module proves uniqueness
for \eqref{eq:motzkin} and identifies that recursion with
\eqref{eq:alternating}. Thus its concluding theorem concerns the original
finite sum. The recurrence assumptions in the general statement specify
only the source sequence. They do not assume any estimate or part of the
conclusion. The separate Challenge module imports only Mathlib. Its
problem also requires existence of a sequence with the source recurrence.
The Solution module proves this problem, including existence and the
sharp threshold.

Table~\ref{tab:lean} gives the principal correspondence. Names are written
relative to the namespace \texttt{MathAndTheory.Motzkin}.

\begin{table}[ht]
\centering\small
\begin{tabular}{>{\raggedright\arraybackslash}p{0.38\linewidth}>{\raggedright\arraybackslash}p{0.54\linewidth}}
\toprule
Mathematical statement & Lean declaration\\\midrule
Source recurrence correspondence &
\path{Source.motzkin_isMotzkin};
\path{Source.eq_motzkin_of_isMotzkin}\\
Source finite-sum correspondence &
\path{Source.alternatingSum_eq_alternating}\\
Strict real-root inequality &
\path{Source.finite_sum_roots_strict_log_concave}\\
Original tail-nine conjecture &
\path{Source.zhao_conjecture_one}\\
Center-five failure &
\path{Source.finite_sum_roots_sharp_failure}\\
Optimal tail start &
\path{Source.finite_sum_tail_start_optimal}\\
Complete challenge &\texttt{Source.solution}\\\bottomrule
\end{tabular}
\caption{Principal formal statements. The source-definition module
passed its native Lean build and axiom audit.}\label{tab:lean}
\end{table}

The underlying modules prove the ratio estimates, positivity, size
bound, and infinite range. They then use proof-producing exact arithmetic
for the remaining indices. The root function is the actual real power
$T_n^{1/n}$. The formal proof does not replace that function with an
uninterpreted symbol or an assumed inequality.

An independent Python program also counts Motzkin paths and applies the
finite alternating sum directly. It agrees with the published OEIS table
through index 300 \cite{OEIS}. It checks the exact inequality at centers
6 through 100 and the reverse comparison at center five. These finite
checks support the result. They do not prove the infinite range.
The Lean statement uses the recurrence characterization. It does not
claim a formal proof of the lattice-path interpretation.

The proof and reproducibility files form one result family.
The supporting ratio lemmas are not separate publication claims.

\section*{Authorship and AI use}

Dev D. Goyal directed the project and is responsible for publication.
The research proofs and original formalizations were produced with
ChatGPT (OpenAI), using the GPT-6 model with Pro reasoning, under his
direction. AI assistance was also used for the subsequent mathematical
audit, the source-definition proof, the independent computation, and this
manuscript. These contributions do not constitute independent human
expert review. The result has not been peer reviewed. No claim of first
priority is made. The repository records the completed review and
verification checks.

\bibliographystyle{plainnat}
\bibliography{references}
\end{document}
